\section{Segunda parte: demostración informal — borrador-esbozo-demostración}

\subsection{Motivación: combinar el razonamiento de alto nivel con la demostración de bajo nivel}

\begin{knowledgebox}{Sledgehammer: herramienta de demostración automática}
Sledgehammer en Isabelle es una herramienta de «martillo» (Hammer) que envía el problema a demostradores de teoremas automáticos externos (lógica de primer orden, lógica de orden superior, resolutores SMT); es muy hábil para llenar pequeños vacíos en una demostración, pero no puede manejar estrategias de demostración de alto nivel que requieren creatividad.
\end{knowledgebox}

\subsection{El método Draft, Sketch, Prove (DSP)}

\begin{importantbox}{Proceso de demostración en tres pasos}
\begin{enumerate}
    \item \textbf{Draft} (borrador): el LLM genera una demostración informal (en lenguaje natural)
    \item \textbf{Sketch} (esbozo): el LLM traduce la demostración informal a un esbozo formal, dejando vacíos por llenar
    \item \textbf{Prove} (demostración): un demostrador de bajo nivel (como Sledgehammer) llena todos los vacíos
\end{enumerate}
Si se llenan todos los vacíos, se obtiene una demostración formal completa.
\end{importantbox}

La estrategia de búsqueda difiere de la búsqueda en árbol tradicional: se generan múltiples borradores informales o múltiples traducciones a esbozos formales, se intenta llenar los vacíos, y basta con que uno tenga éxito para obtener la demostración.

\subsection{Resultados experimentales}

En mini-F2F (en la época de CodeX/Minerva):
\begin{itemize}
    \item El desempeño mejora notablemente conforme aumenta el número de muestreos
    \item Con muestreos suficientes, las demostraciones informales generadas por el LLM incluso superan el desempeño obtenido con demostraciones escritas a mano por humanos
\end{itemize}

